\chapter{Backup, restore és disaster recovery}\label{ch:recovery}

\noindent\textbf{Tanulási út: Gyors út: production alap.} Definiáld a helyreállíthatóságot, mielőtt a production adat pótolhatatlanná válik.\par\medskip

A backup akkor értékes, ha az alkalmazás \textbf{bizonyíthatóan újraindítható egy ismert recovery pontból}. Az üzleti állapot application DB-n, local-auth identity store-on és AppFs-en is átnyúlhat, ezért előre rögzítsd a teljes recovery state-et, a konzisztencia-határt és a restore bizonyításának módját.

\section{Először készíts state inventoryt}

Production előtt írd össze, mi tartozik az alkalmazás állapotához. Tipikusan:

\begin{itemize}
  \item application database;
  \item local-auth SQLite DB, ha local auth használatban van;
  \item AppFs \code{data_root}, például feltöltött média és más runtime adat;
  \item az immutable release azonosítója és artifact hash-e;
  \item a migration set és checksumok;
  \item a trusted config/policy fájlok verziója vagy hash-e;
  \item a backup időpontja, az alkalmazásverzió és a DB backend.
\end{itemize}

A secret fájlok külön secret-management boundaryhoz tartoznak. A recovery runbooknak tudnia kell, honnan állíthatók helyre a szükséges DB-, Redis-, TLS- vagy LDAP-credentialök, de \textbf{secret érték ne kerüljön plaintext release- vagy backup-manifestbe}.

\section{Mi nem üzleti source of truth?}

A V1 Redis-használat session-, rate-limit- és public-cache állapot. Ezeket ne kezeld üzleti backup részeként.

Disaster recovery után biztonságos és elfogadható baseline az üres Redis:

\begin{itemize}
  \item a sessionök megszűnnek, ezért a felhasználók újra belépnek;
  \item a public cache újraépül;
  \item a rate-limit ablakok újraindulnak.
\end{itemize}

Ha ugyanaz a Redis cluster más alkalmazás üzleti adatait is tárolja, annak mentése már annak a rendszernek a saját backup-contractja.

\section{RPO és RTO: a technikai mentés előtt üzleti döntés}

Két számot még a backup implementálása előtt rögzíteni kell:

\begin{description}
  \item[RPO -- Recovery Point Objective] legfeljebb mennyi adatvesztés fogadható el;
  \item[RTO -- Recovery Time Objective] mennyi idő alatt kell helyreállnia a szolgáltatásnak.
\end{description}

Ha például az RPO 24 óra, egy napi backup elvileg elegendő lehet. Ha az RPO néhány perc, már DB-szintű WAL/binlog/PITR vagy más infrastruktúra kellhet. A Velran runtime nem tud szervezeti RPO/RTO-t kitalálni; ezt az operatornak kell a backup technológiára lefordítania.

\section{A V1 biztonságos baseline: kontrollált maintenance backup}

A DB és az AppFs együtt alkothat üzleti állapotot. Egy DB sor például hivatkozhat egy feltöltött képre. Külön időpontban készült DB dump és AppFs tar ezért nem automatikusan konzisztens recovery point.

A legegyszerűbb auditálható V1 baseline:

\begin{lstlisting}
traffic drain
-> SIGTERM / graceful shutdown
-> stopped state verified
-> application DB backup
-> local-auth DB backup, if used
-> AppFs data_root snapshot
-> recovery manifest + hashes
-> service start
-> readiness
\end{lstlisting}

Ez rövid írási szünettel jár, viszont világos konzisztenciahatárt ad. Online backup csak akkor nevezhető konzisztensnek, ha a DB/storage infrastruktúra ugyanahhoz a recovery ponthoz kötött snapshotot biztosít, vagy az alkalmazás adatmodellje bizonyíthatóan tolerálja az eltérő snapshot-időpontokat. A runtime önmagában nem ígér distributed snapshotot.

\section{Application DB: a backend saját backup eszközét használd}

A Velran nem próbálja elrejteni a natív adatbázis-backup mechanizmusokat. Az operator a választott backend támogatott eszközével dolgozzon, lehetőleg külön backup credentialdel.

\subsection*{PostgreSQL}

Egy tipikus logikai backup:

\begin{lstlisting}
pg_dump --format=custom --no-owner --no-acl \
  --file app.dump "$DATABASE_URL"
\end{lstlisting}

Izolált restore-próbához:

\begin{lstlisting}
createdb velran_restore_test
pg_restore --no-owner --no-acl \
  --dbname velran_restore_test app.dump
\end{lstlisting}

A production secretet ne tedd shell historyba. Nagy adatbázisnál fizikai backup és PITR lehet jobb választás; annak runbookja már a PostgreSQL platform része.

\subsection*{MariaDB}

Tranzakcionális InnoDB schema esetén tipikus logikai mentés:

\begin{lstlisting}
mariadb-dump --single-transaction \
  --routines --triggers app > app.sql
\end{lstlisting}

Restore izolált adatbázisba:

\begin{lstlisting}
mariadb velran_restore_test < app.sql
\end{lstlisting}

A \code{--single-transaction} nem tesz nem tranzakcionális táblákat varázsütésre konzisztenssé. Vegyes schema esetén maradjon a kontrollált maintenance baseline.

\subsection*{SQLite}

Élő SQLite adatbázist ne nevezz backupnak csak azért, mert \code{cp}-vel lemásoltad. Használd a SQLite backup mechanizmusát vagy stopped állapotban fájlrendszer-snapshotot.

\begin{lstlisting}
sqlite3 /srv/velran/data/app.db \
  ".backup '/backup/app.db'"
\end{lstlisting}

A restore-próba mindig külön test pathra történjen; ne production DB-n próbáld ki, hogy működik-e a mentés.

\section{A local-auth DB külön security backup}

Local auth használatakor az identity store külön SQLite boundary. Password hash, TOTP secret és recovery hash is lehet benne, ezért backupját \textbf{secretként} kezeld.

\begin{lstlisting}
sqlite3 /srv/velran/auth/users.db \
  ".backup '/backup/local-auth.db'"
\end{lstlisting}

A backup legyen titkosított storage-on, szűk hozzáféréssel és dokumentált retentionnel. Restore után ellenőrizd a fájl és a parent directory ownership/permission értékeit is. Az identity store elvesztése nem rekonstruálható az application DB-ből.

\section{AppFs: a teljes data root tartós állapot lehet}

Ne csak a HTTP-n publikált media alkönyvtárat mentsd. A deklarált AppFs \code{data_root} teljes tartalma lehet recovery state.

Stopped baseline mellett például:

\begin{lstlisting}
tar -C /srv/velran -cpf data-root.tar data
\end{lstlisting}

Restore után őrizd meg vagy állítsd vissza a szükséges ownership/permission értékeket. A service user csak a konfigurációban deklarált writable rootokat írhassa; symlinkkel ne lazítsd az AppFs confinementet.

\section{Minden recovery ponthoz tartozzon manifest}

A backup objektumok önmagukban nem mondják meg, melyik release-hez és migration állapothoz tartoznak. Készíts secretmentes manifestet:

\begin{lstlisting}
backup_id=2026-09-04T180000Z
release_id=2026-09-04.1
server_sha256=...
app_source_sha256=...
migrations_sha256=...
config_sha256=...
db_backend=postgresql
app_db_backup_sha256=...
local_auth_backup_sha256=...
data_root_backup_sha256=...
\end{lstlisting}

A manifest hashokat és azonosítókat tartalmazzon, ne credential-értékeket. A mellékelt \code{deploy/release-record.sh} read-only módon segít a release/source/migration/config állapot rögzítésében.

\section{Restore drill: a backup bizonyítása}

A sikeres backup job még nem bizonyít recovery képességet. Rendszeresen, izolált környezetben restore-próbát kell futtatni.

\begin{lstlisting}
new empty restore DB
-> application DB restore
-> local-auth restore, if used
-> AppFs restore into separate root
-> select exact immutable release
-> restore-test config on isolated host/port
-> velran-server --check-config
-> velran-cli check main.vrn
-> velran-cli migrate verify
-> start velran-server on isolated listener
-> /health/live + /health/ready
-> application-specific smoke test
-> log and audit verification
\end{lstlisting}

A \code{deploy/restore-verify.sh} a nem mutáló config/source/migration ellenőrzések egy részét automatizálja. Nem indít production listenert és nem futtat \code{migrate apply}-t.

A drill eredményét dátummal, backup ID-val és release ID-val dokumentáld. A ``backup elkészült'' nem release gate; a szervezet RPO/RTO elvárásához illő, friss restore-evidence kell.

\section{Upgrade előtt legyen recovery döntés}

Egy kontrollált production upgrade ajánlott sorrendje:

\begin{lstlisting}
immutable release artifact + hashes
-> migration verify
-> compatibility decision
-> consistent backup + recent restore evidence
-> migration apply
-> migration verify
-> atomic release switch
-> controlled restart
-> live + ready + smoke
-> logs + metrics + audit review
\end{lstlisting}

A migration kompatibilitását release-szinten rögzíteni kell. A kulcskérdés: a régi alkalmazás fut-e az új schema mellett?

\section{Expand/contract tartja nyitva a gyors rollback ablakot}

Ha az új schema backward-compatible, egy hibás application release visszakapcsolható a korábbi immutable artifactra DB restore nélkül.

Ajánlott minta:

\begin{enumerate}
  \item \textbf{expand}: új nullable/defaultolt oszlop, tábla vagy index úgy, hogy a régi alkalmazás tovább fusson;
  \item deployold az új alkalmazást;
  \item ha kell, végezz kontrollált adat-backfillt;
  \item csak későbbi release-ben \textbf{contract}: távolítsd el a régi oszlopot/constraintet, amikor a rollback window már lezárult.
\end{enumerate}

Destruktív rename/drop vagy azonnali inkompatibilis shape-váltás bezárhatja az app-only rollback lehetőségét. Ezt release előtt explicit jelöld.

\section{Rollback és restore két külön művelet}

\subsection*{Hibás alkalmazás, backward-compatible schema}

\begin{lstlisting}
traffic drain
-> previous immutable release
-> controlled restart
-> live + ready + smoke
\end{lstlisting}

Ilyenkor nincs DB restore, ezért a hiba után létrejött üzleti adat megmarad.

\subsection*{Hibás alkalmazás, de forward fix lehetséges}

Preferáld az új immutable release-t vagy forward migrationt. Ne köss automatikus reverse SQL-t a release symlink visszaváltásához.

\subsection*{Schema vagy adat sérült}

Ez már disaster recovery:

\begin{lstlisting}
stop writes
-> choose and record recovery point
-> restore DB + auth DB + AppFs consistently
-> select matching release and config
-> verify + smoke
-> controlled return to traffic
\end{lstlisting}

A recovery point után keletkezett írások elveszhetnek. Ez üzleti incidensdöntés, nem technikai részlet: az adatvesztési ablakot explicit vállalni és dokumentálni kell.

\section{Tiltott rövidítések}

\begin{itemize}
  \item élő SQLite DB egyszerű \code{cp} másolatát backupnak nevezni;
  \item csak az application DB-t menteni, ha AppFs is üzleti állapot;
  \item restore-próba nélkül sikeresnek tekinteni a backupot;
  \item automatikus down migrationt kötni application rollbackhez;
  \item production DB-n restore-próbát végezni;
  \item migration credentialt átadni a normál application service-nek;
  \item secret értéket release- vagy backup-manifestbe írni;
  \item Redis session/cache restore-t a disaster recovery feltételévé tenni.
\end{itemize}

\section{Production recovery checklist}

Egy rendszer akkor tekinthető recovery-readynek, ha legalább az alábbiak igazak:

\begin{enumerate}
  \item a state inventory dokumentált;
  \item az RPO és RTO rögzített;
  \item a backup titkosítási, retention- és hozzáférési policy ismert;
  \item van azonosítható, friss sikeres restore drill;
  \item a release migration kompatibilitása dokumentált;
  \item az app-only rollback eljárás próbált;
  \item a destruktív restore tulajdonosa és jóváhagyási útja ismert;
  \item Redis elvesztésének hatása elfogadott;
  \item a recovery runbookból nem hiányzik a release/config/secret visszaállítási forrás;
  \item incidens után live, ready, smoke, log és audit ellenőrzés is történik.
\end{enumerate}

A recovery legfontosabb elve: \textbf{a mentés nem fájl, hanem bizonyított visszaállítási képesség}. Ha a restore drill nem fut, az RPO/RTO csak feltételezés.

\section{Gyakorlat: restore-t bizonyíts, ne backupot}
\paragraph{Gyakorlási mód: üzemeltetési szcenárió.} Használd az előszóban meghatározott üzemeltetési módszert.

Állíts vissza egy alkalmazásállapot-másolatot izolált környezetbe, majd ellenőrizd a rekordszámokat, kritikus üzleti invariánsokat, media referenciákat és az alkalmazás indulását. Mérd a restore időt, és jegyezd fel az incidenskezelést blokkoló manuális lépéseket.

\section{Tipikus hiba: komponenseket mentünk közös recovery modell nélkül}
Adatbázissorok, feltöltött fájlok, konfiguráció és release identity együtt alkothatnak egy helyreállítható állapotot akkor is, ha külön tárolódnak. A backup terv mondja meg, hogyan illeszkednek össze restore idején.

\section{Folyamatos mintaprojekt: Secure Notes}
Végezz teljes Secure Notes recovery próbát a tartós adattal és az esetleges media attachmentekkel együtt. Ellenőrizd, hogy a helyreállított generáció ugyanazokat az ownership és publication szabályokat szolgálja ki, mint a szimulált hiba előtt.
